\chapter{모스 부호: \nt{GLUT} 뉴런과 \nt{GABA} 뉴런이 말하는 언어에 대해 우리가 아는 것}
\chaptermark{모스 부호}
\label{sec:code}

\section{기호 하나짜리 알파벳}

모스 부호에는 점과 대시라는 두 기호가 있고, 그 사이에 세 가지 길이의 쉼이 있다. \ref{sec:neurontypes}장에서 보았듯이 뉴런의 알파벳은 더 빈약하다. 기호가 하나뿐이다. 스파이크는 약 1밀리초 동안 지속되고, 한 뉴런의 스파이크는 높이와 모양이 모두 같다. 그러므로 메시지 전체는 쉼에, 즉 시각의 열 $t_1,t_2,t_3,\dots$에 담긴다. 대시가 없고 점의 리듬만이 의미를 싣는 부호이다(그림~\ref{fig:morse}).

\begin{figure}[t]
\centering
\begin{tikzpicture}[>=Stealth,x=0.08cm,y=1cm]
% Morse letter: dot, dash, dot
\node[left,font=\small] at (-6,2.55) {모스:};
\fill[black] (0,2.45) rectangle (6,2.65);
\fill[black] (14,2.45) rectangle (32,2.65);
\fill[black] (40,2.45) rectangle (46,2.65);
\node[right,font=\small,gray!40!black] at (52,2.55) {점, 대시, 점: 의미는 길이와 쉼에 있다};
% stimulus onset
\node[left,font=\small] at (-6,0.4) {뉴런:};
\draw[->,thick,green!45!black] (0,-0.55) -- (0,-0.05);
\node[below,font=\scriptsize,green!45!black] at (0,-0.55) {자극};
\draw[gray!70!black] (0,0) -- (152,0) node[right,black] {\small $t$};
% spikes: single ones blue, the burst red
\foreach \s in {14,26,41,88,112,131}{\draw[very thick,blue!60!black] (\s,0) -- (\s,0.8);}
\foreach \s in {60,63,66,69}{\draw[very thick,red!70!black] (\s,0) -- (\s,0.8);}
% latency
\draw[<->,thick] (0,1.1) -- node[above,font=\scriptsize] {$t_1$: 첫 스파이크 시각} (14,1.1);
% burst
\draw[decorate,decoration={brace,amplitude=4pt},red!70!black] (58,0.95) -- (71,0.95) node[midway,above=4pt,font=\scriptsize] {버스트: `대시'};
% counting windows
\foreach \a/\b/\n in {0/50/3,50/100/5,100/150/2}{
  \draw[decorate,decoration={brace,amplitude=4pt,mirror},gray!50!black] (\a+1,-0.2) -- (\b-1,-0.2) node[midway,below=4pt,font=\small] {$n=\n$};}
\node[font=\scriptsize,gray!40!black] at (75,-1.05) {발화율 부호: 창마다 스파이크 수 $n$};
\foreach \x in {0,50,100,150}{\draw[gray!60!black] (\x,-0.06) -- (\x,0.06);}
\end{tikzpicture}
\caption{모스 부호와 뉴런의 부호. 같은 스파이크 열을 여러 방식으로 읽을 수 있다. $50\ms$ 창마다 스파이크를 세거나(\ref{sec:rate}절), 지연 $t_1$을 재거나~\eqref{eq:latency}, `대시'인 버스트를 알아챌 수 있다~\eqref{eq:burst}.}
\label{fig:morse}
\end{figure}

어떤 쉼이 가능한가? 아래로는 불응기가 제한한다. 스파이크 뒤 약 1밀리초 동안 뉴런은 다시 발화할 수 없으므로 초당 수백 개보다 많은 스파이크는 나오지 않는다. 위로는 아무 제한이 없다. 뉴런은 몇 분 동안 침묵할 수 있다. 피질 \nt{GLUT} 뉴런의 발화율은 초당 1개 미만에서 수십 개 사이에 있고, 대부분의 뉴런은 대부분의 시간을 아래쪽 경계 근처에서 작동한다.

\ref{sec:glutcomputer}장과 \ref{sec:gaba}장에서 우리는 스파이크로 수를 적는 방식을 그때그때 하나씩 가정했다. 이제 직접 묻자. \nt{GLUT} 뉴런과 \nt{GABA} 뉴런은 서로 어떤 언어로 말하는가? 완전한 답은 없다. 이 언어는 해독되지 않았다. 그러나 확실히 알려진 것도 있으며, 알려진 것의 거의 전부는 통신 엔지니어처럼 추론해서 얻을 수 있다. 채널 용량, 잡음, 수신기, 전송 비용이다.

\section{얼마나 전달할 수 있는가}

상한부터 시작하자. 수신자가 스파이크 시각을 정밀도 $\Delta t$로 구별한다고 하자. $\Delta t$보다 작은 이동은 수신자에게 구별되지 않는다. 시간을 불응기보다 짧은 길이 $\Delta t$의 칸으로 자르면, 각 칸에는 스파이크가 하나 있거나 없다(그림~\ref{fig:cells}). 평균 발화율이 $r$이면 스파이크가 칸에 들어갈 확률은 $p=r\Delta t$이다. 흐름이 가장 많은 정보를 나르는 것은 칸들이 독립일 때이며, 이때 각 칸은 $p\ll1$에서 엔트로피 $-p\log_2p-(1-p)\log_2(1-p)\approx p\log_2(e/p)$비트를 준다. $\Delta t$로 나누면 전송 속도의 한계와 스파이크 하나당 몫을 얻는다~\cite{MacKay1952}.
\begin{empheq}[box=\keybox]{equation}
R_{\max} \;\approx\; r\,\log_2\frac{e}{r\,\Delta t}\ \ \text{bit/s},
\qquad
\frac{R_{\max}}{r} \;\approx\; \log_2\frac{e}{r\,\Delta t}\ \ \text{스파이크당 bit}.
\label{eq:spikeentropy}
\end{empheq}
$r$이 초당 스파이크 10개이고 $\Delta t=1\ms$이면, 스파이크당 약 8비트, 초당 약 80비트이다.

\begin{figure}[t]
\centering
\begin{tikzpicture}[>=Stealth,x=0.5cm,y=1cm]
% time cut into cells of width Delta t; a cell holds one impulse or none
\foreach \k in {0,...,23}{\draw[gray!55] (\k,0) rectangle (\k+1,0.7);}
\foreach \k in {2,7,8,15,21}{
  \fill[red!10] (\k,0) rectangle (\k+1,0.7);
  \draw[very thick,red!70!black] (\k+0.5,0.05) -- (\k+0.5,0.65);}
\foreach \k in {0,...,23}{
  \pgfmathtruncatemacro{\bit}{(\k==2)||(\k==7)||(\k==8)||(\k==15)||(\k==21)}
  \ifnum\bit=1 \def\bitcol{red!70!black}\else \def\bitcol{gray!60!black}\fi
  \node[font=\scriptsize,text=\bitcol] at (\k+0.5,-0.25) {\bit};}
\draw[<->,gray!60!black] (0,0.95) -- (1,0.95) node[midway,above,font=\scriptsize,black] {$\Delta t$};
\node[font=\small,anchor=east] at (-0.3,0.35) {스파이크};
\node[font=\small,anchor=east] at (-0.3,-0.25) {비트};
\draw[decorate,decoration={brace,amplitude=5pt,mirror},gray!60!black] (0,-0.55) -- (24,-0.55)
  node[midway,below=6pt,font=\scriptsize,black,align=center] {세기만 하는 수신자: `$n=5$';\\ 칸을 보는 수신자: $24$칸 중 정확히 어느 $5$칸인가, 즉 $\log_2\binom{24}{5}\approx15$비트};
\end{tikzpicture}
\caption{용량~\eqref{eq:spikeentropy}은 어디서 오는가. 시간을 길이 $\Delta t$의 칸으로 자르면, 칸마다 스파이크가 있거나($1$) 없다($0$). 스파이크 흐름은 이진 문자열이다. 스파이크가 칸에 들어갈 확률은 $p=r\Delta t$이다. 칸이 잘수록 문자열이 길어지고 비트가 많아진다. 스파이크 계수기는 1이 \emph{어디에} 있는지에 적힌 것을 거의 모두 잃는다.}
\label{fig:cells}
\end{figure}

스파이크당 비트 수는 시간 정밀도에 로그로만 의존한다. 정밀도가 $10\ms$이면 스파이크당 약 5비트가 남는다. 그러나 수신자가 1초 동안의 스파이크를 세기만 한다면, $r=10$에서 1초 전체에 대해 2비트도 얻지 못한다(\ref{sec:rate}절).

\begin{keyblock}
\begin{proposition}[스파이크 채널의 용량]
\label{prop:capacity}
평균 발화율 $r$인 스파이크 흐름을 시간 정밀도 $\Delta t$로 읽으면 \eqref{eq:spikeentropy}보다 많이 나르지 못한다. 이 용량의 거의 전부는 스파이크의 \emph{시각}에 들어 있다. 스파이크를 세기만 하는 수신자는 같은 흐름에서 시각을 1밀리초 정밀도로 구별하는 수신자보다 수십 배 적게 뽑아낸다.
\end{proposition}
\end{keyblock}

이것은 상한이지 측정이 아니다. 자극을 알고 있는 감각 뉴런에서 측정하면 스파이크당 1비트에서 몇 비트가 나온다. 가장 좋은 경우 그 정밀도에 대해 계산한 한계의 절반까지 이른다~\cite{Rieke1997}. 따라서 적어도 신경계의 입구에서는 스파이크 시각이 버려지지 않고 쓰인다.

\section{잡음 있는 채널}

용량~\eqref{eq:spikeentropy}은 잡음 없이 계산했다. 잡음이 어디서 생기는지에 대해서는 세 가지 사실이 확실히 알려져 있다.

\emph{스파이크 발생기 자체는 정확하다.} 조직 조각과 함께 뇌에서 꺼낸 피질 뉴런에 시간에 따라 빠르게 변하는 같은 전류를 여러 번 주입하면, 스파이크는 약 1밀리초의 정밀도로 반복된다~\cite{Mainen1995}. 전류가 일정하면 스파이크 시각은 시행마다 흩어진다. 정확한 시각을 만드는 것은 뉴런 자체가 아니라 입력의 빠른 변화이다.

\emph{시냅스는 믿을 수 없다.} \nt{GLUT} 시냅스에 도착한 스파이크는 어떤 확률 $p$로만 신경전달물질 한 꾸러미를 내보낸다. 해마 시냅스에서는 스파이크의 절반 넘게가 수신자에게 그냥 전달되지 않으며, 원인은 전도 실패가 아니라 바로 방출의 무작위성이다~\cite{AllenStevens1994}. 통신 엔지니어에게 이것은 소거 채널이다. 두 뉴런 사이의 여러 시냅스가 사정을 어느 정도 구해 준다.

\emph{살아 있는 피질에서 스파이크 흐름은 무작위로 보인다.} 스파이크 간격은 무작위 푸아송 흐름과 거의 같게 흩어져 있고, 같은 자극을 되풀이하면 창 안의 스파이크 수가 분산이 평균에 가깝게 변한다~\cite{Softky1993}.

정확한 소자가 어떻게 무작위 흐름을 낳는가? 뉴런은 입력 수천 개를 받고, 하나하나는 믿을 수 없으며, 흥분과 억제는 \ref{sec:gaba}장에서처럼 균형을 이룬다. 막전압은 문턱값 근처에서 흔들리고, 스파이크 시각은 이 흔들림이 정한다. 이것이 잡음인지 우리가 읽지 못하는 메시지인지는 대체로 열린 문제이다. `잡음'의 일부는 이미 메시지로 밝혀졌다. 시각 피질에서 변동의 상당 부분은 동물 자신의 움직임을 따라간다~\cite{Stringer2019}. 여기에는 원리적인 어려움이 있다. 잘 압축된 메시지는 부호를 모르는 사람에게 무작위 흐름과 구별되지 않는다.

\section{발화율: 느리지만 믿을 만하다}
\label{sec:rate}

가장 단순한 부호는 \emph{발화율 부호}이다. 수는 뉴런이 창 $T$ 동안 내는 스파이크의 개수에 적힌다. 소거도 시각의 흔들림도 이런 부호에는 두렵지 않다. 그러나 대가가 있다. 흐름이 푸아송이면 창 안의 스파이크 수 $n$은 평균 $rT$, 산포 $\sqrt{rT}$를 가진다.
\begin{empheq}[box=\keybox]{equation}
\frac{\sigma_n}{\bar n} \;=\; \frac{1}{\sqrt{r\,T}} .
\label{eq:rateerror}
\end{empheq}
발화율을 10퍼센트 정밀도로 알려면 스파이크 백 개가 필요하다. 초당 스파이크 20개라면 5초이다. 이웃한 수준은 산포 두 배쯤 떨어져 있어야 구별되므로, 발화율 $r$까지 시간 $T$ 동안 약 $\sqrt{rT}$개의 수준이 구별된다. $r=10$, $T=1\,\text{s}$이면 서너 수준, 2비트 미만이다.
뇌는 그렇게 느리게 작동하지 않는다. 사람은 한순간 보여 준 그림에서 동물을 알아보며, 뇌 속의 응답은 약 $150\ms$ 뒤에 나타난다~\cite{Thorpe1996}. 거칠게 추정하면 그동안 신호는 연속된 처리 단계 열 개쯤을 한 단계에 $10$--$20\ms$씩 지나간다. 피질의 보통 발화율에서 각 뉴런은 한 단계 동안 스파이크를 0개나 1개 낼 시간밖에 없고, 이런 창에서 그 발화율은 정의되지 않는다. 출구는 두 가지이다. 뉴런을 많이 쓰거나, 시각을 보는 것이다.

\emph{많은 뉴런.} 뉴런 $N$개가 같은 수를 나르고 수신자가 그 스파이크를 더한다고 하자. 잡음이 독립이면 \eqref{eq:rateerror}에서 $rT$ 대신 $NrT$가 들어간다. $r=20$인 뉴런 백 개는 $15\ms$ 동안 스파이크 서른 개를 주어 약 20퍼센트의 정밀도를 낸다. 공간이 시간을 대신한다. 그러나 이웃한 피질 뉴런의 잡음은 독립이 아니다. 이웃한 쌍의 스파이크 수 상관 계수는 보통 약 $0.1$이다~\cite{Zohary1994}. 상관 계수가 $c$이면 뉴런 $N$개에 대한 평균의 분산은 다음과 같다.
\begin{empheq}[box=\keybox]{equation}
\sigma^2_N \;=\; \sigma^2_1\,\frac{1+(N-1)\,c}{N}
\;\xrightarrow[N\to\infty]{}\; c\,\sigma^2_1 .
\label{eq:correlated}
\end{empheq}

\begin{keyblock}
\begin{proposition}[평균의 한계]
\label{prop:pooling}
잡음이 상관되어 있으면 무리에 대한 평균은 오차를 기껏해야 $1/\sqrt{c}$배 줄인다. $c=0.1$이면 세 배이며, 이 이득은 이미 뉴런 수십 개에서 달성된다. 더 보태 봐야 소용없다.
\end{proposition}
\end{keyblock}

모든 상관이 똑같이 해로운 것은 아니다. 정밀도를 제한하는 것은 공통 잡음 가운데 \emph{신호 자체를 닮은} 부분, 즉 측정하는 양이 변했을 때와 같은 방식으로 무리 전체를 미는 부분뿐이다~\cite{MorenoBote2014}. 뇌에 이런 잡음이 실제로 얼마나 있는지는 아직 밝히는 중이다.

많은 뉴런을 쓰는 다른 방법도 있다. 소리 방향 검출기(그림~\ref{fig:itd})처럼 수를 \emph{어느} 뉴런이 발화했는지에 적는 것이다. 이것이 \emph{위치} 부호이며, 알려진 것 중 가장 믿을 만하다. 감각 뉴런에서 전선의 번호는 피부의 어느 점이 닿았는지, 소리의 높이가 얼마인지를 말해 주며, 이는 여러 번 확인되었다. 뉴런 수로 보면 비싸지만 빠르다. 메시지 하나에 지연 하나가 든다.

\section{시각: 빠르지만, 어디서부터 재는가}

두 번째 출구는 수를 스파이크 시각에 적는 것이다. 뉴런~\eqref{eq:glutneuron}에 시각 $t=0$부터 전압을 수준 $U$까지 올릴 일정한 입력을 가하자. 전압은 $V=U\bigl(1-e^{-t/\tau}\bigr)$로 커지고 다음 시각에 문턱값 $\theta$에 이른다.
\begin{empheq}[box=\keybox]{equation}
t_1 \;=\; \tau\,\ln\frac{U}{U-\theta}, \qquad U>\theta .
\label{eq:latency}
\end{empheq}
강한 입력은 이른 스파이크, 약한 입력은 늦은 스파이크, 문턱 아래 입력은 스파이크 없음이다. $\tau=10\ms$에서 입력 $U=4\theta$는 $2.9\ms$ 뒤에, $U=2\theta$는 $6.9\ms$ 뒤에, $U=1.2\theta$는 $17.9\ms$ 뒤에 첫 스파이크를 준다. 스파이크 단 하나가 수를 나른다.

그런데 $t_1$을 어느 순간부터 재는가? 공통 시계는 없으며, 수신자는 자극이 언제 시작되었는지 모른다. 답은 세 가지이다.

\emph{상대 지연.} 두 뉴런이 같은 자극을 보면, 두 지연의 차 $t_1^A-t_1^B$는 시작 시각이 아니라 입력에만 의존한다. 시각은 지연이 있는 동시성 검출기가 비교한다(그림~\ref{fig:itd}). 뉴런 $N$개가 스파이크를 하나씩 냈다면, 그 도착 \emph{순서}만으로 최대 $\log_2N!$비트를 나른다. 뉴런 열 개면 약 22비트인데, `발화했는가 아닌가'라는 답은 10비트를 넘지 못한다. 망막에서는 출력 뉴런들의 첫 스파이크 상대 지연으로 그림을 빠르고 믿을 만하게 복원할 수 있음이 밝혀졌다~\cite{Gollisch2008}.

\emph{낱말의 시작을 알리는 일제 발화.} 자극이 급격히 바뀌면 수많은 뉴런이 거의 동시에 발화한다. 이런 일제 발화는 모스 부호에서 낱말 사이의 긴 쉼처럼 그 자체로 시작 표시가 되고, 수신자는 거기서부터 지연을 잰다.

\emph{국소 리듬.} \ref{sec:gaba}장에서 \nt{GLUT}--\nt{GABA} 루프는 주기 $12$--$50\ms$의 국소 리듬을 낳았다. 이것은 클록이 아니지만, 그 한 주기 안에서 입력이 강한 뉴런은 입력이 약한 뉴런보다 먼저 발화하며, \eqref{eq:latency}는 \emph{위상} 부호가 된다. 수는 스파이크가 주기의 어느 부분에 왔는지에 적힌다. 이런 부호는 관찰되었다. 공간의 특정 장소를 맡은 뉴런은 동물이 `자기' 장소를 지나가면서 느린 국소 리듬의 주기 안에서 점점 더 일찍 발화한다~\cite{OKeefe1993}. 뇌가 위상을 얼마나 널리 쓰는지는 아직 모른다.

\section{부호는 수신자가 고른다}

스파이크 열에는 발화율도, 시각도, 순서도 있다. 그중 무엇을 읽을지는 수신자가 정하며, 파라미터 하나, 즉 자신의 시간 상수 $\tau$로 정한다.

방정식~\eqref{eq:glutneuron}를 시간에 대해 평균하자. 발화율 $r_j$인 독립 흐름들이 뉴런에 들어오면 평균 전압과 그 상대적 흔들림은 다음과 같다.
\begin{empheq}[box=\keybox]{equation}
\bar V \;=\; \tau\sum_j w_j\,r_j ,
\qquad
\frac{\sigma_V}{\bar V} \;\approx\; \frac{1}{\sqrt{2\,r_{\text{in}}\,\tau}} ,
\label{eq:reader}
\end{empheq}
여기서 두 번째 식은 가중치가 같을 때의 것이고, $r_{\text{in}}$은 모든 입력의 발화율 합이다. $\tau$가 크면 전압은 거의 흔들리지 않고 발화율의 가중합을 따라간다. 수신자는 \emph{적분기}이며, 발화율을 읽고 시각은 잊는다. $\tau$가 작으면 전압이 스파이크 사이에 떨어질 시간이 있어서, 여러 스파이크가 창~\eqref{eq:window} 안에 들어왔을 때만 문턱값에 이른다. 수신자는 \emph{동시성 검출기}이며, 시각을 읽는다(그림~\ref{fig:reader}). 초당 스파이크 5개인 입력 천 개는 $\tau=10\ms$에서 약 10퍼센트의 흔들림을 주며, 거의 순수한 적분이다. \ref{sec:gaba}장에서처럼 억제가 $\tau$를 $2\ms$로 줄이면 흔들림은 20퍼센트 남짓으로 커지고, 뉴런은 일치를 알아채기 시작한다.

\begin{figure}[t]
\centering
\begin{tikzpicture}[>=Stealth,x=0.115cm,y=1cm]
% common input: the same spike train for both receivers
\foreach \s in {6,21,22,38,52,55.5,59,62.5,66,69.5,73,76.5,80,83.5,87,90.5,94,97.5}{\draw[thick,gray!40!black] (\s,5.25) -- (\s,5.65);}
\draw[gray!70!black] (0,5.25) -- (105,5.25);
\node[left,font=\small] at (-1,5.45) {입력};
\node[above,font=\scriptsize,gray!40!black] at (13.5,5.65) {드묾};
\node[above,font=\scriptsize,gray!40!black] at (21.5,6.0) {쌍};
\node[above,font=\scriptsize,gray!40!black] at (75,5.65) {잦음};
% axes and thresholds
\foreach \y/\th in {2.7/2.5,0/1.5}{
  \draw[gray!70!black] (0,\y) -- (106,\y) node[right,black] {\small $t$};
  \draw[densely dashed,gray!60!black] (0,\y+0.6*\th) -- (104,\y+0.6*\th) node[right,black,font=\small] {$\theta$};
}
\node[anchor=south west,font=\small] at (0,4.3) {$\tau=10\ms$: 적분기};
\node[anchor=south east,font=\small] at (104,1.0) {$\tau=2\ms$: 동시성 검출기};
% integrator: tau=10 ms, theta=2.5w
\begin{scope}[yshift=2.7cm]
\foreach \a/\b/\v in {0/6/0.000,6/21/1.000,21/22/1.223,22/38/2.107,38/52/1.425,52/55.5/1.351,55.5/59/1.952,59/62.5/2.376,62.5/66/0.000,66/69.5/1.000,69.5/73/1.705,73/76.5/2.201,76.5/80/0.000,80/83.5/1.000,83.5/87/1.705,87/90.5/2.201,90.5/94/0.000,94/97.5/1.000,97.5/104/1.705}{\draw[very thick,blue!60!black] plot[domain=\a:\b,samples=14] (\x,{0.6*\v*exp(-(\x-\a)/10)});}
\foreach \s/\p/\q in {6/0.000/1.000,21/0.223/1.223,22/1.107/2.107,38/0.425/1.425,52/0.351/1.351,55.5/0.952/1.952,59/1.376/2.376,62.5/1.674/2.500,62.5/2.500/0.000,66/0.000/1.000,69.5/0.705/1.705,73/1.201/2.201,76.5/1.551/2.500,76.5/2.500/0.000,80/0.000/1.000,83.5/0.705/1.705,87/1.201/2.201,90.5/1.551/2.500,90.5/2.500/0.000,94/0.000/1.000,97.5/0.705/1.705}{\draw[very thick,blue!60!black] (\s,0.6*\p) -- (\s,0.6*\q);}
\foreach \s in {62.5,76.5,90.5}{\draw[->,thick,red!70!black] (\s,1.58) -- (\s,1.95);}
\end{scope}
% coincidence: tau=2 ms, theta=1.5w
\begin{scope}[yshift=0cm]
\foreach \a/\b/\v in {0/6/0.000,6/21/1.000,21/22/1.001,22/38/0.000,38/52/1.000,52/55.5/1.001,55.5/59/1.174,59/62.5/1.204,62.5/66/1.209,66/69.5/1.210,69.5/73/1.210,73/76.5/1.210,76.5/80/1.210,80/83.5/1.210,83.5/87/1.210,87/90.5/1.210,90.5/94/1.210,94/97.5/1.210,97.5/104/1.210}{\draw[very thick,blue!60!black] plot[domain=\a:\b,samples=14] (\x,{0.6*\v*exp(-(\x-\a)/2)});}
\foreach \s/\p/\q in {6/0.000/1.000,21/0.001/1.001,22/0.607/1.500,22/1.500/0.000,38/0.000/1.000,52/0.001/1.001,55.5/0.174/1.174,59/0.204/1.204,62.5/0.209/1.209,66/0.210/1.210,69.5/0.210/1.210,73/0.210/1.210,76.5/0.210/1.210,80/0.210/1.210,83.5/0.210/1.210,87/0.210/1.210,90.5/0.210/1.210,94/0.210/1.210,97.5/0.210/1.210}{\draw[very thick,blue!60!black] (\s,0.6*\p) -- (\s,0.6*\q);}
\foreach \s in {22}{\draw[->,thick,red!70!black] (\s,0.98) -- (\s,1.35);}
\end{scope}
\foreach \x in {0,50,100}{\draw[gray!60!black] (\x,-0.06) -- (\x,0.06) node[below,font=\scriptsize] {\x\,ms};}
\end{tikzpicture}
\caption{같은 입력, 서로 다른 두 수신자. 스파이크마다 전압이 $w$만큼 올랐다가 뉴런 모델~\eqref{eq:glutneuron}에서처럼 흘러내린다. 빨간 화살표는 수신자의 스파이크이다. 느린 수신자($\tau=10\ms$, $\theta=2.5w$)는 스파이크가 \emph{많은} 곳에서 발화하고 쌍은 알아채지 못한다. 빠른 수신자($\tau=2\ms$, $\theta=1.5w$)는 창~\eqref{eq:window} 안의 쌍에만 발화하고, 잦지만 일치하지 않는 흐름은 흘려보낸다.}
\label{fig:reader}
\end{figure}

측정에 따르면 활성 피질에서 막 컨덕턴스는 휴지 상태보다 몇 배 커진다~\cite{Destexhe2003}. 그러니 거기서는 $\tau$가 더 짧고, 작동 중인 피질 뉴런은 적분기보다 동시성 검출기에 더 가깝다.

\begin{keyblock}
\begin{proposition}[부호는 수신자가 정한다]
\label{prop:reader}
같은 스파이크 열이 느린 수신자에게는 발화율을, 빠른 수신자에게는 시각을, 신경전달물질이 방출되는 조직 공간에게는 몇 초에 걸쳐 평활된 농도 수준을 나른다(\ref{sec:serocomputer}장). 어떤 부호가 쓰이는지는 송신자가 아니라 수신자의 시간 상수가 정하며, 그 시간 상수는 억제가 조정한다.
\end{proposition}
\end{keyblock}
\section{점과 대시: 필터로서의 시냅스}

지금까지 시냅스는 가중치가 달린 전선이었다. 실제로는 시냅스의 응답이 앞선 스파이크에 의존하며, 이것이 뉴런의 언어에 두 번째 기호를 준다.

\emph{고갈: 시냅스는 변화를 전달한다.} 방출할 때마다 방출 준비가 된 신경전달물질 저장량의 몫 $U$가 쓰이고~\cite{Tsodyks1997}, 저장량은 수백 밀리초 정도의 시간 상수 $\tau_{\text{rec}}$로 회복된다. 발화율이 $r$이면 스파이크당 응답 진폭은 대략 $A(r)=A_0/(1+U r\tau_{\text{rec}})$ 수준에 자리 잡으며, 여기서 $A_0$는 쉬고 난 시냅스의 응답이다. 수신자가 단위 시간에 받는 전류는 다음과 같다.
\begin{empheq}[box=\keybox]{equation}
r\,A(r) \;=\; \frac{A_0\,r}{1+U r\,\tau_{\text{rec}}}
\;\xrightarrow[r\to\infty]{}\; \frac{A_0}{U\tau_{\text{rec}}} .
\label{eq:depression}
\end{empheq}
$U=0.5$, $\tau_{\text{rec}}=0.8\,\text{s}$이면 포화는 이미 초당 스파이크 $1/(U\tau_{\text{rec}})=2.5$개에서 시작된다. 발화율이 5에서 20으로 오르면 정상 상태 전류는 겨우 3분의 1 늘어난다. 대신 새 발화율의 첫 스파이크들은 저장량이 아직 예전 그대로일 때 도착하므로 전류가 한순간 네 배로 뛰었다가 $\tau_{\text{rec}}/(1+Ur\tau_{\text{rec}})\approx90\ms$ 만에 줄어든다(그림~\ref{fig:depression}).

이런 시냅스는 발화율이 아니라 그 \emph{변화}, 본질적으로 상대 변화 $\Delta r/r$를 전달한다~\cite{Abbott1997depression}. 전자공학자에게 이것은 전선 안에 바로 넣은 고역 통과 필터이다.

\begin{figure}[t]
\centering
\begin{tikzpicture}[>=Stealth,x=12.5cm,y=1cm]
% input rate
\draw[->,gray!70!black] (0,2.6) -- (0.9,2.6) node[right,black] {\small $t$};
\draw[->,gray!70!black] (0,2.6) -- (0,4.3) node[above,black] {\small $r$};
\draw[very thick,blue!60!black] (0,2.9) -- (0.2,2.9) -- (0.2,3.8) -- (0.8,3.8);
\node[left,font=\scriptsize] at (0,2.9) {$5$};
\node[left,font=\scriptsize] at (0,3.8) {$20$};
\node[right,font=\small,blue!60!black] at (0.4,4.1) {시냅스 입력의 발화율};
% output current
\draw[->,gray!70!black] (0,0) -- (0.9,0) node[right,black] {\small $t$};
\draw[->,gray!70!black] (0,0) -- (0,2.45) node[left,black] {\small $rA$};
\draw[densely dashed,gray!60!black] (0,0.667) -- (0.8,0.667);
\draw[very thick,red!70!black] (0,0.5) -- (0.2,0.5) -- (0.2,2.0)
   plot[domain=0.2:0.8,samples=60] (\x,{0.667+(2.0-0.667)*exp(-(\x-0.2)/0.089)});
\node[left,font=\scriptsize] at (0,0.5) {$1.7$};
\node[left,font=\scriptsize] at (0,2.0) {$6.7$};
\node[right,font=\scriptsize] at (0.8,0.667) {$2.2$};
\node[right,font=\small,red!70!black] at (0.28,1.6) {수신자로 가는 전류: 도약 뒤 $\approx90\ms$ 만에 감소};
\draw[gray!60!black] (0.2,-0.08) -- (0.2,0.08); \node[below,font=\scriptsize] at (0.2,0) {$0.2\,\text{s}$};
\draw[gray!60!black] (0.8,-0.08) -- (0.8,0.08); \node[below,font=\scriptsize] at (0.8,0) {$0.8\,\text{s}$};
\end{tikzpicture}
\caption{$U=0.5$, $\tau_{\text{rec}}=0.8\,\text{s}$일 때 식~\eqref{eq:depression}에 따른 고갈 시냅스. 전류는 초당 $A_0$ 단위이고, 파선은 정상 상태 전류이다. 정상 상태 전류는 $1.7$에서 $2.2$로 늘었을 뿐이지만, 도약하는 순간에는 $6.7$까지 오른다. 시냅스는 수신자에게 발화율이 얼마인지가 아니라 발화율이 바뀌었다는 것을 알린다.}
\label{fig:depression}
\end{figure}

\emph{촉진: 대시로서의 버스트.} 다른 시냅스들은 방출 확률이 낮지만, 스파이크마다 수십 밀리초 동안 그 확률이 커진다. 이런 시냅스를 지나는 단일 스파이크는 대개 사라진다. 그러나 \emph{버스트}, 즉 몇 밀리초 간격으로 연이어 나오는 스파이크 몇 개는 거의 확실히 통과한다. 촉진이 없어도 버스트의 스파이크 $k$개가 모두 사라질 확률은
\begin{empheq}[box=\keybox]{equation}
P_{\text{loss}} \;=\; (1-p)^k ,
\label{eq:burst}
\end{empheq}
이며, $p=0.2$이면 단일 스파이크에서 $0.8$, 네 개짜리 버스트에서 $0.41$이다. 촉진은 이 확률을 더 줄인다. 이렇게 뉴런에게 두 번째 기호가 생긴다. 단일 스파이크는 점, 버스트는 대시이다. 고갈 시냅스는 쉼 뒤의 첫 스파이크를 가장 잘 전달하고, 촉진 시냅스는 버스트를 통과시키고 홀로 된 점은 죽인다. 버스트가 더 믿을 만하게 전달된다는 것은 측정되었다. 뇌가 실제로 버스트를 별도의 기호로 쓴다는 것은 그럴듯한 가설이다~\cite{Lisman1997}.

\section{\nt{GABA} 뉴런은 어떻게 말하는가}

피질에서 가장 수가 많은 것은 빠른 \nt{GABA} 뉴런이다~\cite{Tremblay2016}. 이들의 스파이크는 \nt{GLUT} 뉴런보다 짧고, 초당 수백 개의 발화율로 고르게, 오랫동안, 거의 지치지 않고 스파이크를 낼 수 있다. 이들의 시냅스는 더 믿을 만하다. 각 수신자와의 연결이 많은 방출 자리로 이루어져 있어 스파이크가 사라지는 일이 거의 없다. 이들의 출력은 수신자에서 스파이크가 생겨나는 자리 가까이에 붙으므로, 수신자가 입력을 어떻게 더하는지가 아니라 발화할지, 언제 발화할지를 제어한다.

이들 자신은 이웃한 많은 \nt{GLUT} 뉴런에게서 흥분을 받고 주변의 총활동을 알린다. \ref{sec:gaba}장의 나눗셈에 딱 필요한 것이다. 끝으로 빠른 \nt{GABA} 뉴런은 서로 연결되어 쉽게 함께 발화한다. 앞에서 말한 국소 리듬을 정하는 것이 바로 이들이다~\cite{Cardin2009}.

모스 부호의 언어로 말하면 \nt{GABA} 뉴런은 글자가 아니라 \emph{쉼}을 맡는다. 이들은 \nt{GLUT} 스파이크의 끊임없는 흐름을 수신자가 발화할 수 있는 창들로 자르고, 일치 창을 좁히며, 흐름이 너무 시끄러워지면 전체를 낮춘다.

\begin{keyblock}
\begin{proposition}[역할 분담]
\label{prop:punctuation}
\nt{GLUT} 뉴런은 메시지의 내용, 즉 \emph{무엇}과 \emph{얼마나}를 나른다. 빠른 \nt{GABA} 뉴런은 메시지의 표지, 즉 \emph{언제} 말해도 되는지와 \emph{얼마나 크게}를 나른다. 또 다른 부류는 억제하는 뉴런을 억제하여 관문이 \emph{누구에게} 열리는지를 정한다. 이 분담의 개별 부분은 측정되었다. 일반 규칙으로서는 작업 가설이다.
\end{proposition}
\end{keyblock}

더 느린 \nt{GABA} 뉴런들도 있는데, 그 출력은 수신자의 개별 입력에 붙는다. 이들은 \ref{sec:gaba}장의 금지~\eqref{eq:veto}처럼 작동한다. \nt{GLUT} 스파이크의 흐름 하나를 닫고 나머지는 건드리지 않는다.

\nt{GABA} 뉴런을 빠른 것과 느린 것으로 나누는 것은 오래된 그림이며, 완전하지 않다. 지금은 피질에서 적어도 세 개의 큰 부류를 구별하며~\cite{Tremblay2016}, 세 번째 부류는 뜻밖의 구조를 가진다. 이 부류는 \nt{GLUT} 뉴런이 아니라 다른 \nt{GABA} 뉴런, 무엇보다 입력을 닫는 뉴런들을 억제한다~\cite{Pfeffer2013}. 억제의 억제는 이중 부정이다. 이런 뉴런은 수신자에게 흥분성 스파이크를 하나도 보내지 않고 관문을 \emph{연다}. 이 뉴런을 켜는 것은 피질의 다른 영역에서 오는 신호와 브로드캐스트 채널이므로, 이 뉴런은 네트워크 한 구역 전체의 허가 입력처럼 작동한다. 이것이 활성인 동안 금지가 풀리고 흐름이 통과한다~\cite{Pi2013}. 부록~\ref{sec:othermediators}에서는 이 기법을 탈억제라고 부른다.
\section{절약하는 언어}

뉴런의 언어에 대해 확실히 알려진 마지막 사실은 그것이 비싸다는 것이다. 스파이크 하나와 그것이 일으키는 시냅스의 일에는 ATP 분자가 약 $10^9$개 들고(설치류 피질에 대한 추정~\cite{Attwell2001}), 분자 하나는 약 $10^{-19}$줄을 낸다. 따라서 스파이크 하나의 값은 $E_{\text{spk}}\sim10^{-10}$줄 정도이다. 뇌 전체는 약 $20$와트를 쓰며, 그 절반이 신호 전달에 쓰인다면, 즉 $P\sim10$와트이면 뉴런 하나의 평균 발화율은 다음과 같다.
\begin{empheq}[box=\keybox]{equation}
\bar r \;\approx\; \frac{P}{N\,E_{\text{spk}}}
\;\sim\; \frac{10\ \text{W}}{8.6\cdot10^{10}\cdot10^{-10}\ \text{J}}
\;\approx\; 1\ \text{spike/s}.
\label{eq:energy}
\end{empheq}
피질에 대한 더 자세한 추정은 이보다도 적은 값을 준다~\cite{Lennie2003}. 뇌의 부호는 \emph{희소}할 수밖에 없다. 뉴런 가운데 작은 몫만이 빠르게 작동할 수 있다.

\eqref{eq:spikeentropy}를 보면 이것이 그리 나쁘지 않다. 정밀도가 $1\ms$일 때, 초당 $100$번 발화하는 뉴런의 스파이크는 5비트를 넘지 못하지만, 초당 한 번 발화하는 뉴런의 스파이크는 11비트까지 나른다. 스파이크마다 값을 치러야 한다면 드물게 말하는 편이 이득이다. 피질은 실제로 정밀도를 에너지와 맞바꾼다. 먹이가 부족하면 뉴런은 스파이크 발화율은 유지하되 시냅스 전류에 덜 쓰고, 그 응답은 덜 정밀해진다~\cite{Padamsey2022}.

모스 부호에서 자주 쓰는 글자는 짧다. 영어 글에서 가장 흔한 글자 E는 점 하나이다. 알려진 바로는 뉴런의 부호도 다른 형태로 같은 규칙을 따른다. 예상된 것에는 스파이크를 적게 쓴다~\cite{Barlow1961}. 자극이 일정하면 뉴런은 발화율을 점점 낮추고, \ref{sec:glutcomputer}장의 센서는 수준이 아니라 변화에 응답하며, \ref{sec:neurontypes}장의 \nt{DOPA} 뉴런은 보상 예측 오차만을 알린다.

\begin{keyblock}
\begin{proposition}[절약]
\label{prop:economy}
에너지는 뉴런의 평균 발화율을 초당 스파이크 1개 정도로 제한한다. 그래서 \nt{GLUT} 뉴런의 언어는 희소하며, 스파이크를 새 소식, 즉 변한 것이나 예측되지 않은 것에 쓴다. 에너지 한계는 측정되었다. 뇌의 부호가 줄당 비트 수를 최대로 하도록 맞추어져 있다는 것은 근거가 탄탄한 가설이다.
\end{proposition}
\end{keyblock}

\section{요약}

\begin{table}[t]
\centering
\footnotesize
\setlength{\tabcolsep}{4pt}
\renewcommand{\arraystretch}{1.15}
\begin{tabular}{>{\raggedright}p{2.6cm}>{\raggedright}p{2.3cm}>{\raggedright}p{3.0cm}>{\raggedright}p{2.0cm}>{\raggedright\arraybackslash}p{3.6cm}}
\toprule
적는 방식 & 나르는 것 & 읽는 쪽 & 메시지당 시간 & 알려진 것 \\
\midrule
발화한 뉴런의 번호(위치 부호) & 어떤 값인가 & 모든 수신자; 동시성 검출기 & 지연 하나 & 감각 뉴런과 소리 방향 검출에서 확립됨 \\
뉴런 하나의 발화율 & 크기 & $\tau$가 큰 적분기; 조직 공간(\ref{sec:serocomputer}장) & 수백 ms--수 초 & 확립됨; $10$--$20\ms$의 처리 단계에는 너무 느림 \\
무리의 발화율 & 크기 & 입력이 수천 개인 뉴런 & $10$--$50\ms$ & 확립됨; 이득은 상관~\eqref{eq:correlated}으로 제한됨 \\
첫 스파이크 시각, 순서 & 크기, 순서 & 지연이 있는 동시성 검출기 & 지연 하나 & 신경계 입구에서 확립됨; 피질에서는 논의 중 \\
국소 리듬 속의 위상 & 크기, 위치 & 공통 리듬을 가진 뉴런 & 주기 $12$--$50\ms$ 및 더 느림 & 일부 영역에서 관찰됨; 일반적 역할은 가설 \\
버스트(`대시') & `중요함': 믿을 만한 전달 & 촉진 시냅스 & $\sim10\ms$ & 신뢰성은 측정됨; 별도 기호라는 것은 가설 \\
발화율의 변화 & 새 소식 & 고갈 시냅스 & $\sim0.1\,\text{s}$ & 확립됨 \\
빠른 \nt{GABA} 뉴런의 스파이크 & 언제, 얼마나 크게 & 이웃한 모든 \nt{GLUT} 뉴런 & $1$--$30\ms$ & 리듬과 나눗셈은 측정됨; 규칙으로서의 `구두점'은 가설 \\
\bottomrule
\end{tabular}
\caption{\nt{GLUT} 뉴런과 \nt{GABA} 뉴런이 스파이크로 메시지를 적는 방식. 오른쪽 열은 측정된 것과 추정된 것을 구분한다.}
\label{tab:code}
\end{table}

표~\ref{tab:code}는 우리가 아는 것을 모은 것이며, 우리가 모르는 것도 거기서 보인다. 피질이 무리의 발화율뿐 아니라 정확한 스파이크 시각을 얼마나 쓰는지 우리는 모른다. 뉴런에게 `낱말', 즉 전체로 읽히는 여러 뉴런 스파이크의 안정된 조합이 있는지 모른다. 그리고 뇌의 여러 부분에서 언어가 같은지도 모른다. 모스 부호는 그 표가 알려져 있기에 하룻저녁이면 익힐 수 있다. 뉴런 언어의 표는 아직 만들어야 하며, 그것을 만드는 사람은 아마 통신 엔지니어일 것이다.

\section{연습문제}

\begin{exercise}[\lvlA]
\label{ex:04:1}
\emph{숫자로 본 용량.} $\Delta t=1\ms$. (a)~스파이크 값이 \eqref{eq:energy}와 같을 때, 발화율이 초당 스파이크 $1$개인 뉴런은 줄당 몇 비트를 주는가? (b)~$r=10$에서 1초 동안의 스파이크를 세기만 하는 수신자는 한계~\eqref{eq:spikeentropy}보다 몇 배 적게 뽑아내는가?
\end{exercise}

\begin{exercise}[\lvlB]
\label{ex:04:2}
\emph{최적 발화율.} 뉴런은 전력 $P_0+rE_{\text{spk}}$를 쓴다. 여기서 $P_0$는 $20$~W의 나머지 절반을 모든 뉴런에 나눈 값이고, $E_{\text{spk}}=10^{-10}$~J이다. $\Delta t=1\ms$인 \eqref{eq:spikeentropy}에 따른 줄당 비트 수 $R_{\max}(r)/(P_0+rE_{\text{spk}})$는 어떤 발화율 $r$에서 가장 큰가?
\end{exercise}

\begin{exercise}[\lvlB]
\label{ex:04:3}
\emph{어떤 상관이 해로운가.} 뉴런 $N=1000$개, 분산 $\sigma^2=1$, 쌍별 상관 $c=0.1$. 기울기가 모두 같은 $f'=\mathbf 1$일 때와, $\mathbf 1^{\mathsf T}f'=0$인 기울기 $\pm1$일 때의 피셔 정보 $\mathcal I=f'^{\mathsf T}\Sigma^{-1}f'$($\Sigma$는 공분산 행렬)을 계산하라. 둘은 몇 배 다른가?
\end{exercise}

\begin{exercise}[\lvlB]
\label{ex:04:4}
\emph{모델링: 동시성 검출기.} 뉴런~\eqref{eq:glutneuron}가 초당 스파이크 $5$개인 푸아송 입력 $1000$개를 받고 $w=1$이다. 문턱값 $\theta$는 자유 전압이 1초에 한 번 넘도록 잡는다. $\tau=10$과 $2\ms$에서 동시 입력 몇 개 $m$이 확률 $1/2$로 뉴런을 발화시키는지 모델로 구하라.
\end{exercise}

\begin{exercise}[\lvlB]
\label{ex:04:5}
\emph{설계: 상대 변화 검출기.} 시냅스~\eqref{eq:depression}를 고르라. 초당 스파이크 $40$개에서의 정상 상태 전류가 $10$개에서보다 $10\,\%$ 넘게 높지 않아야 하고, $40$으로 뛰어오른 뒤 전류가 시간 상수 $50\ms$ 이하로 새 수준에 이르러야 한다. 가장 작은 $\tau_{\text{rec}}$는 얼마이고, 그때 $U$는 얼마인가?
\end{exercise}

\begin{exercise}[\lvlA]
\label{ex:04:6}
\emph{첫 스파이크 시각과 버스트.} (a)~\eqref{eq:latency}에서 첫 스파이크가 정확히 시간 상수 하나 뒤에 오게 하는 입력을 구하라. 이 입력은 무엇에 의존하는가? (b)~방출 확률이 $p=0.2$일 때, 모두 잃을 확률이 $5\,\%$ 아래가 되려면 버스트에 스파이크가 몇 개 있어야 하는가?
\end{exercise}

\begin{exercise}[\lvlC]
\label{ex:04:7}
\emph{탐구: 배치에는 몇 비트가 있는가.} 뉴런은 독립 칸들~\eqref{eq:spikeentropy}이고, $r=10$, $\Delta t=1\ms$이다. (a)~칸 $20$개로 된 `낱말'의 엔트로피를 스파이크 수의 엔트로피와 나머지로 나누어라. (b)~기록 $N=30$개와 $10^4$개로부터 낱말 빈도로 엔트로피를 추정하는 것을 시뮬레이션하라. 추정치는 얼마나 낮게 나오는가?
\end{exercise}
